% Data sources and the twelve data contracts of docs/contracts/catalog.md.

\section{Data and data contracts}\label{sec:data}

\subsection{Data sources and splits}\label{sec:data-sources}

DinHot reads no external dataset. Its inputs are the \nCases{} case records of the case library, written to a source document by the pipeline, and the published anchors that the acceptance criteria quote. The owner of the case library is the product itself (the repository named in the header block, licence Apache-2.0); each record states its parameters in the units of Section~\ref{sec:notation}, and every record is synthetic, so no sensitive or measured sample is involved. The records span \minTraps{} to \maxTraps{} traps and four categories of cases, listed in Section~\ref{sec:protocol}. The disordered array of case C13 reproduces the geometry of Kim et al.~\cite{kim2019} with the coordinates drawn by DinHot, because the published ones are not available. The quasi-periodic set of the benchmark is likewise DinHot's own, since the 119 traps of Polin et al.~\cite{polin} are not published as coordinates.

The derived data are the committed artifacts of the build. The directory of derived data holds \derivedFiles{} files of \derivedKiB{} KiB in all, the largest trace being \traceMaxKiB{} KiB against the lane gate of 256 KiB, and the released graph of the learned model is \cghBytes{} MB. The artifacts are written by the pipeline and read by the web, the API and the tests; they are committed so that a reader can re-check every number without running the bake, and the pipeline tests run in a sandbox that fails the test run if a committed artifact changes.

The splits matter only for the two learned models. The training data are generated by a seeded generator in the data root and never committed. The held-out sets are written apart before training, with seeds other than the training ones, and their hash is recorded in the model card: \cghN{} volumes for the three-plane network and \wsN{} geometries for the warm start. A model is judged on the held-out set only, and it is released only if it meets its criterion; otherwise the baseline answers behind the same contract. The acceptance tests of the other methods use no split: a test fixes a task and a seed, and the benchmark repeats the tasks over up to \benchSeeds{} seeds, with the tuning of a method run on other seeds (offset by a thousand) at 30 iterations, so that no tuned row is scored on the seeds that tuned it.

Each data contract of the catalog is described below in the order of the catalog: purpose, fields (type, unit, range, in the table of the contract; the nested fields of the larger contracts are in Appendix~\ref{app:fields}), producer and consumers, version and the behaviour when data are missing. A schema file is the single source of its contract: the Python producers and consumers validate against it in their tests, and the TypeScript types are generated from it. A change bumps the contract's version by the semantic-version rule (a removed or renamed field, a narrowed type or a new required field is a major step, a new optional field a minor step, a description change a patch) and updates both sides in one change. Table~\ref{tab:contracts} lists the twelve contracts.

\begin{table}[t]\centering\footnotesize
\caption{The data contracts of the build, from the catalog: identifier, version, owner (the module that produces the document or checks the boundary) and the committed documents the schema governs.}\label{tab:contracts}
\begin{tabular}{@{}l l >{\raggedright\arraybackslash}p{0.34\textwidth} >{\raggedright\arraybackslash}p{0.34\textwidth}@{}}\toprule Contract & Version & Owner & Documents governed\\\midrule
\texttt{params-record} & 1.0.0 & {\scriptsize\path{data-pipeline/pipeline/io/contract.py}} & {\scriptsize\path{run-time exchange, no document}}\\
\texttt{case-index} & 1.0.0 & {\scriptsize\path{data-pipeline/pipeline/core/manifest.py}} & {\scriptsize\path{data/derived/manifests/index.json}}\\
\texttt{case-manifest} & 1.0.0 & {\scriptsize\path{data-pipeline/pipeline/core/manifest.py}} & {\scriptsize\path{data/derived/manifests/*.json; !data/derived/manifests/index.json}}\\
\texttt{trace} & 1.0.0 & {\scriptsize\path{data-pipeline/pipeline/core/trace.py}} & {\scriptsize\path{data/derived/*/trace.json}}\\
\texttt{benchmark-results} & 1.0.0 & {\scriptsize\path{data-pipeline/pipeline/stages/benchmark.py}} & {\scriptsize\path{data/derived/benchmark/results.json}}\\
\texttt{cgh-net} & 1.0.1 & {\scriptsize\path{data-pipeline/pipeline/stages/train.py}} & {\scriptsize\path{models/cgh-net.json; models/cgh-net.onnx}}\\
\texttt{warmstart-model} & 1.0.1 & {\scriptsize\path{data-pipeline/pipeline/stages/train.py}} & {\scriptsize\path{models/warmstart.json; models/warmstart.onnx}}\\
\texttt{model-metrics} & 1.0.1 & {\scriptsize\path{data-pipeline/pipeline/stages/evaluate.py}} & {\scriptsize\path{data/derived/models/metrics.json}}\\
\texttt{model-evidence} & 1.0.0 & {\scriptsize\path{data-pipeline/pipeline/stages/train.py}} & {\scriptsize\path{data/derived/models/*-training.json; data/derived/models/*-rejected.json; data/derived/models/evaluation-detail.json}}\\
\texttt{engine-reference} & 1.0.0 & {\scriptsize\path{data-pipeline/pipeline/stages/engine_reference.py}} & {\scriptsize\path{tests/fixtures/engine/*.json}}\\
\texttt{readiness} & 1.0.0 & {\scriptsize\path{data-pipeline/pipeline/core/readiness.py}} & {\scriptsize\path{data/derived/readiness.json}}\\
\texttt{api} & 1.0.0 & {\scriptsize\path{app/}} & {\scriptsize\path{docs/contracts/examples/api-unavailable.json#UnavailableBody}}\\
\bottomrule\end{tabular}\end{table}


\subsection{Case record (params-record)}\label{sec:data-params-record}

The params-record contract (version \cver{params-record}) is the ingestion boundary: the record of one case, with the optical setup, the device, the target (traps or a shaped potential), the aberration, the method set and budgets, the seed, the feedback, dynamics and rearrangement blocks and the expected ranges of the metrics. Its fields are in Table~\ref{tab:contract-params-record}, and the nested ones in Appendix~\ref{app:fields}. The producer is the module that validates records against the schema with a standard-library validator, and the consumers are the bake and the tests. The contract states rules a schema cannot: a trap at or beyond the half field is refused with the field's bounds, which case C11 exists to prove. The missing-data behaviour is partial: a rejected record is dropped with its reason and the bake goes on with the accepted records, and the readiness document lists the rejected case ids; if the source document of cases is absent, the bake uses the registry's cases and readiness says missing.

\begin{table}[t]\centering\footnotesize
\caption{Top-level fields of the contract \texttt{params-record}, read from its schema file: type, requirement, unit (a dash marks an identifier, a count or a dimensionless quantity) and the range or constant the schema states.}\label{tab:contract-params-record}
\begin{tabular}{@{}l l c l >{\raggedright\arraybackslash}p{0.3\textwidth}@{}}\toprule Field & Type & Req. & Unit & Range\\\midrule
\texttt{\scriptsize case\_id} & string & yes & - & pattern\\
\texttt{\scriptsize setup} & object & yes & - & -\\
\texttt{\scriptsize device} & object or null & yes & - & -\\
\texttt{\scriptsize target} & object & yes & - & -\\
\texttt{\scriptsize aberration} & object or null & yes & - & -\\
\texttt{\scriptsize methods} & array of string & yes & - & 1 to 24 items\\
\texttt{\scriptsize reference\_method} & string & yes & - & pattern\\
\texttt{\scriptsize budget} & object & yes & - & -\\
\texttt{\scriptsize seed} & integer & yes & - & 0 to 2147483647\\
\texttt{\scriptsize feedback} & object or null & yes & - & -\\
\texttt{\scriptsize dynamics} & object or null & yes & - & -\\
\texttt{\scriptsize rearrangement} & object or null & yes & - & -\\
\texttt{\scriptsize expect} & object & yes & - & -\\
\bottomrule\end{tabular}\end{table}


\subsection{Case index (case-index)}\label{sec:data-case-index}

The case-index contract (version \cver{case-index}) is the inventory of the build: one entry per case, with bilingual titles and categories, the default case, and the status (baked or refused) with the reason for a refused one. Its fields are in Table~\ref{tab:contract-case-index}. The producer is the manifest module of the pipeline and the consumer is the frontend, which reads the index first, so that without it there is no site. The behaviour when it is missing is therefore the strictest of the catalog: the site build stops at the copy of the data, the API answers 503 with the state missing, and readiness says missing. When the index is stale, because its engine version is older than the product's, readiness says stale with both versions and the views show it. A refused case has an entry and no manifest, and the workbench shows its reason.

\begin{table}[t]\centering\footnotesize
\caption{Top-level fields of the contract \texttt{case-index}, read from its schema file: type, requirement, unit (a dash marks an identifier, a count or a dimensionless quantity) and the range or constant the schema states.}\label{tab:contract-case-index}
\begin{tabular}{@{}l l c l >{\raggedright\arraybackslash}p{0.3\textwidth}@{}}\toprule Field & Type & Req. & Unit & Range\\\midrule
\texttt{\scriptsize schema} & const & yes & - & = dinhot.index/v1\\
\texttt{\scriptsize engine\_version} & string & yes & - & pattern\\
\texttt{\scriptsize default\_case} & string & yes & - & length >= 1\\
\texttt{\scriptsize cases} & array of object & yes & - & 1 to items\\
\bottomrule\end{tabular}\end{table}


\subsection{Case manifest (case-manifest)}\label{sec:data-case-manifest}

The case-manifest contract (version \cver{case-manifest}) describes one baked case: its parameters and seed, the engine version, the lane label and the gate that produced it, the score per method (the metrics scored on the case's device), the models behind the numbers, the expected ranges and the bilingual band, and the published anchor where one exists. The fields are in Table~\ref{tab:contract-case-manifest}, with the nested ones in Appendix~\ref{app:fields}. The manifest carries no clock: the run time that decides the lane is stored as a lane label and as the deterministic budgets, never as milliseconds, so that two bakes of the same inputs write identical bytes, which a test compares byte by byte. If a manifest that the index lists is absent, the contract is partial: readiness lists its path, the case selector marks the case as not available yet and the other cases work, and the API answers 404 for it. A stale manifest is handled as for the index.

\begin{table}[t]\centering\footnotesize
\caption{Top-level fields of the contract \texttt{case-manifest}, read from its schema file: type, requirement, unit (a dash marks an identifier, a count or a dimensionless quantity) and the range or constant the schema states.}\label{tab:contract-case-manifest}
\begin{tabular}{@{}l l c l >{\raggedright\arraybackslash}p{0.3\textwidth}@{}}\toprule Field & Type & Req. & Unit & Range\\\midrule
\texttt{\scriptsize schema} & const & yes & - & = dinhot.manifest/v1\\
\texttt{\scriptsize case\_id} & string & yes & - & length >= 1\\
\texttt{\scriptsize title} & object & yes & - & -\\
\texttt{\scriptsize category} & object & yes & - & -\\
\texttt{\scriptsize real\_or\_synthetic} & enum & yes & - & \{real, synthetic\}\\
\texttt{\scriptsize expected\_band} & object & yes & - & -\\
\texttt{\scriptsize engine} & object & yes & - & -\\
\texttt{\scriptsize params} & object & yes & - & -\\
\texttt{\scriptsize seed} & integer & yes & - & -\\
\texttt{\scriptsize artifact} & object & yes & - & -\\
\texttt{\scriptsize lane} & enum & yes & - & \{live, precompute\}\\
\texttt{\scriptsize gate} & object & yes & - & -\\
\texttt{\scriptsize flags} & array of object & yes & - & -\\
\texttt{\scriptsize methods} & object & yes & - & -\\
\texttt{\scriptsize models} & object & yes & - & -\\
\texttt{\scriptsize expect} & object & yes & - & -\\
\texttt{\scriptsize anchor} & object & yes & - & -\\
\bottomrule\end{tabular}\end{table}


\subsection{Trace (trace)}\label{sec:data-trace}

The trace contract (version \cver{trace}) is the compact replay artifact of one case, at most 256 KiB. It holds, per method, the metric series against iteration (at most 200 points, first and last kept), the final per-trap intensities and metrics; for the case's reference method, the quantized phase mask (256 levels, base64), the focal intensity in decibels relative to the mean trap on a crop of at most 128 bins, with blocks max-pooled so that a peak is never lost, and the axial cut for three-dimensional cases; and, for the dynamic cases, the bead or atom trajectory (at most 4000 points), the power spectrum (at most 200 log-spaced bins) and the rearrangement frames. Table~\ref{tab:contract-trace} lists the top-level fields. A trace that would exceed the byte gate is squeezed through a fixed list of five settings, and one that still does not fit raises an error, so the bake fails and the lane gate never sees an oversized artifact. When an artifact that a manifest names is absent, the case's replay views show not available yet with the reason and the live lane still runs from the manifest's parameters.

\begin{table}[t]\centering\footnotesize
\caption{Top-level fields of the contract \texttt{trace}, read from its schema file: type, requirement, unit (a dash marks an identifier, a count or a dimensionless quantity) and the range or constant the schema states.}\label{tab:contract-trace}
\begin{tabular}{@{}l l c l >{\raggedright\arraybackslash}p{0.3\textwidth}@{}}\toprule Field & Type & Req. & Unit & Range\\\midrule
\texttt{\scriptsize schema} & const & yes & - & = dinhot.trace/v1\\
\texttt{\scriptsize case\_id} & string & yes & - & length >= 1\\
\texttt{\scriptsize grid} & object & yes & - & -\\
\texttt{\scriptsize traps} & array of object & yes & - & 0 to 4096 items\\
\texttt{\scriptsize shape} & object or null & yes & - & -\\
\texttt{\scriptsize methods} & object & yes & - & -\\
\texttt{\scriptsize reference} & object & yes & - & -\\
\texttt{\scriptsize feedback} & object or null & yes & - & -\\
\texttt{\scriptsize dynamics} & object or null & yes & - & -\\
\texttt{\scriptsize rearrangement} & object or null & yes & - & -\\
\texttt{\scriptsize summary} & object & yes & - & -\\
\bottomrule\end{tabular}\end{table}


\subsection{Benchmark results (benchmark-results)}\label{sec:data-benchmark-results}

The benchmark-results contract (version \cver{benchmark-results}) holds the table of the benchmark protocol: one row per method, geometry, device, budget kind, budget and parameter set, with the median and the inter-quartile range over the seeds of the efficiency, the uniformity, the relative standard deviation, the wall-clock time and the iterations, the number of seeds of the cell, the stored Pareto fronts and the cells that did not run, with the reason (Table~\ref{tab:contract-benchmark-results}). The producer is the benchmark stage and the consumer is the benchmark views of the frontend. A cell that cannot run is listed in the skipped list and is absent from the rows, never silently: readiness lists the cells as method and geometry, and the views mark each one as not run. If the document is absent, the views say not available yet with the reason and never show a number from an older bake.

\begin{table}[t]\centering\footnotesize
\caption{Top-level fields of the contract \texttt{benchmark-results}, read from its schema file: type, requirement, unit (a dash marks an identifier, a count or a dimensionless quantity) and the range or constant the schema states.}\label{tab:contract-benchmark-results}
\begin{tabular}{@{}l l c l >{\raggedright\arraybackslash}p{0.3\textwidth}@{}}\toprule Field & Type & Req. & Unit & Range\\\midrule
\texttt{\scriptsize schema} & const & yes & - & = dinhot.benchmark/v1\\
\texttt{\scriptsize engine\_version} & string & yes & - & pattern\\
\texttt{\scriptsize protocol} & object & yes & - & -\\
\texttt{\scriptsize geometries} & array of object & yes & - & -\\
\texttt{\scriptsize rows} & array of object & yes & - & -\\
\texttt{\scriptsize pareto} & array of object & yes & - & -\\
\texttt{\scriptsize skipped} & array of object & no & - & -\\
\bottomrule\end{tabular}\end{table}


\subsection{Model cards (cgh-net and warmstart-model)}\label{sec:data-cgh-net}

Two model slots share the model-card schema (Table~\ref{tab:contract-cgh-net}; the warm-start card is in Table~\ref{tab:contract-warmstart-model}). The cgh-net contract (version \cver{cgh-net}) maps the target amplitudes of three focal planes at $-8$, $0$ and $+8$ SLM-limited Rayleigh ranges, unit-sum intensities, for a $128\times128$ SLM with no padding, to one phase in radians; the baseline behind it is the volumetric Gaussian GS at 30 iterations. The card names the ONNX graph, its size, its sha256 and its opset, the grid, the input and output tensors with their dtype and shape, the training record and the held-out set. The behaviour when the slot has no valid released model (the card or the graph is absent or fails its schema or its checksum) is the state training with the baseline provenance: the baseline answers, and the views label the result with the model that answered. A marker file means that the released model answers while a newer one trains.

\begin{table}[t]\centering\footnotesize
\caption{Top-level fields of the contract \texttt{cgh-net}, read from its schema file: type, requirement, unit (a dash marks an identifier, a count or a dimensionless quantity) and the range or constant the schema states.}\label{tab:contract-cgh-net}
\begin{tabular}{@{}l l c l >{\raggedright\arraybackslash}p{0.3\textwidth}@{}}\toprule Field & Type & Req. & Unit & Range\\\midrule
\texttt{\scriptsize schema} & const & yes & - & = dinhot.model-card/v1\\
\texttt{\scriptsize slot} & enum & yes & - & \{cgh-net, warmstart-model\}\\
\texttt{\scriptsize onnx} & string & yes & - & pattern\\
\texttt{\scriptsize sha256} & string & yes & - & pattern\\
\texttt{\scriptsize bytes} & integer & yes & - & 1 to\\
\texttt{\scriptsize opset} & integer & yes & - & 1 to\\
\texttt{\scriptsize inputs} & array of object & yes & - & 1 to items\\
\texttt{\scriptsize outputs} & array of object & yes & - & 1 to items\\
\texttt{\scriptsize grid} & const & no & - & = 128\\
\texttt{\scriptsize trained\_on} & object & yes & - & -\\
\texttt{\scriptsize held\_out} & object & yes & - & -\\
\texttt{\scriptsize metrics\_path} & string & yes & - & length >= 1\\
\bottomrule\end{tabular}\end{table}


\subsection{Warm-start model (warmstart-model)}\label{sec:data-warmstart-model}

The warmstart-model contract (version \cver{warmstart-model}) maps a padded list of up to 1024 traps (position in focal bins and in Rayleigh ranges, plus the weight) and a mask to one phase per trap in radians, with the random trap phases of the portable generator as its baseline. It shares the card schema of the previous contract and its missing-data behaviour, with the baseline being the random start. In this build the slot is refuted and no model file is committed: the record of the refusal holds the held-out metrics and the control, and the contract test asserts that no warm-start model file exists, so the state of the slot is the baseline answering. The contract is nevertheless documented because the schema, the example and the swap test make the refusal a tested path and not an omission.

\begin{table}[t]\centering\footnotesize
\caption{Top-level fields of the contract \texttt{warmstart-model}, read from its schema file: type, requirement, unit (a dash marks an identifier, a count or a dimensionless quantity) and the range or constant the schema states.}\label{tab:contract-warmstart-model}
\begin{tabular}{@{}l l c l >{\raggedright\arraybackslash}p{0.3\textwidth}@{}}\toprule Field & Type & Req. & Unit & Range\\\midrule
\texttt{\scriptsize schema} & const & yes & - & = dinhot.model-card/v1\\
\texttt{\scriptsize slot} & enum & yes & - & \{cgh-net, warmstart-model\}\\
\texttt{\scriptsize onnx} & string & yes & - & pattern\\
\texttt{\scriptsize sha256} & string & yes & - & pattern\\
\texttt{\scriptsize bytes} & integer & yes & - & 1 to\\
\texttt{\scriptsize opset} & integer & yes & - & 1 to\\
\texttt{\scriptsize inputs} & array of object & yes & - & 1 to items\\
\texttt{\scriptsize outputs} & array of object & yes & - & 1 to items\\
\texttt{\scriptsize grid} & const & no & - & = 128\\
\texttt{\scriptsize trained\_on} & object & yes & - & -\\
\texttt{\scriptsize held\_out} & object & yes & - & -\\
\texttt{\scriptsize metrics\_path} & string & yes & - & length >= 1\\
\bottomrule\end{tabular}\end{table}


\subsection{Held-out metrics (model-metrics)}\label{sec:data-model-metrics}

The model-metrics contract (version \cver{model-metrics}) holds the held-out evaluation of the model slots against their baselines: the number of held-out cases, the accuracy ratio against the baseline at equal wall-clock, the time of the model and the hash of the held-out file, for the slots that have a released model (Table~\ref{tab:contract-model-metrics}). The producer is the evaluate stage; the consumer is the benchmark views of the frontend. A slot with no released model has no entry and its views say so, because a metric for a model that is not released would be a number about a model nobody can use. The comparison it records is always against the baseline of the slot at equal wall-clock, never against a fixed number. If the document is absent, the views say not available yet and never show a number from an older bake; if its engine version is older than the product's, readiness says stale.

\begin{table}[t]\centering\footnotesize
\caption{Top-level fields of the contract \texttt{model-metrics}, read from its schema file: type, requirement, unit (a dash marks an identifier, a count or a dimensionless quantity) and the range or constant the schema states.}\label{tab:contract-model-metrics}
\begin{tabular}{@{}l l c l >{\raggedright\arraybackslash}p{0.3\textwidth}@{}}\toprule Field & Type & Req. & Unit & Range\\\midrule
\texttt{\scriptsize schema} & const & yes & - & = dinhot.model-metrics/v1\\
\texttt{\scriptsize engine\_version} & string & yes & - & pattern\\
\texttt{\scriptsize slots} & object & yes & - & -\\
\bottomrule\end{tabular}\end{table}


\subsection{Training evidence (model-evidence)}\label{sec:data-model-evidence}

The model-evidence contract (version \cver{model-evidence}) is the evidence behind the metrics: the training records of the released model, the records of the rejected candidates (the single-plane network and the warm start) and the evaluation detail with the ratio by kind of target (Table~\ref{tab:contract-model-evidence}). The producer is the training stage and the consumers are the tests of the two learned methods, which read the records to assert the criteria: the committed metrics of the volume network, and the measured negative of the warm start together with its control. A slot with no training record has no evidence to show, and the benchmark page says that the slot answers with its baseline, as for the metrics document. A rejected candidate is kept on purpose: a negative result is a result, and the record lets a reader re-check the refusal.

\begin{table}[t]\centering\footnotesize
\caption{Top-level fields of the contract \texttt{model-evidence}, read from its schema file: type, requirement, unit (a dash marks an identifier, a count or a dimensionless quantity) and the range or constant the schema states.}\label{tab:contract-model-evidence}
\begin{tabular}{@{}l l c l >{\raggedright\arraybackslash}p{0.3\textwidth}@{}}\toprule Field & Type & Req. & Unit & Range\\\midrule
\bottomrule\end{tabular}\end{table}


\subsection{Engine reference values (engine-reference)}\label{sec:data-engine-reference}

The engine-reference contract (version \cver{engine-reference}) holds the reference values that the Python engine writes and the TypeScript parity tests read: for each quantity, the inputs, the expected outputs and the tolerance, absolute and relative (Table~\ref{tab:contract-engine-reference}). It is a test input and never a run-time document, so its missing-data behaviour is that readiness says missing when the folder of fixtures holds no file and the parity tests fail naming the absent file. The fixtures are the means by which the live engine is held to the reference: \liveFixtures{} reference documents are committed, with a loosest absolute tolerance of $\liveTol$, and \liveTests{} browser-side tests compare the port with them (Section~\ref{sec:validation}). The documents are written by the Python engine from fixed inputs and seeds, so that a change of the engine that moves a result beyond the stated tolerance fails a test in the repository instead of surfacing in the browser. The version of the contract is that of the schema, and a change in a tolerance is a change in the contract.

\begin{table}[t]\centering\footnotesize
\caption{Top-level fields of the contract \texttt{engine-reference}, read from its schema file: type, requirement, unit (a dash marks an identifier, a count or a dimensionless quantity) and the range or constant the schema states.}\label{tab:contract-engine-reference}
\begin{tabular}{@{}l l c l >{\raggedright\arraybackslash}p{0.3\textwidth}@{}}\toprule Field & Type & Req. & Unit & Range\\\midrule
\texttt{\scriptsize schema} & const & yes & - & = dinhot.engine-ref/v1\\
\texttt{\scriptsize engine\_version} & string & yes & - & pattern\\
\texttt{\scriptsize function} & string & yes & - & length >= 1\\
\texttt{\scriptsize tolerance} & object & yes & - & -\\
\texttt{\scriptsize cases} & array of object & yes & - & 1 to items\\
\bottomrule\end{tabular}\end{table}


\subsection{Readiness (readiness)}\label{sec:data-readiness}

The readiness contract (version \cver{readiness}) says what is real in a build, row by row, in one of six declared states: available, partial, stale, provider, training and missing (Table~\ref{tab:contract-readiness}). The document is computed from the catalog and the repository, and the front end and the API show it; editing the catalog makes it stale until it is recomputed, which a guard fails on. Of the \nContracts{} rows it judges, \nContractsAvailable{} are fully available in this build, and the build declares the state of the others: \readinessOthers. If the document is missing, the site build stops, and if it cannot be loaded at run time the views keep working and each says readiness unknown. A provider row, a row served by a stand-in while the real producer is not built, is labelled on every surface and never counted as done; in this build every row has a real producer.

\begin{table}[t]\centering\footnotesize
\caption{Top-level fields of the contract \texttt{readiness}, read from its schema file: type, requirement, unit (a dash marks an identifier, a count or a dimensionless quantity) and the range or constant the schema states.}\label{tab:contract-readiness}
\begin{tabular}{@{}l l c l >{\raggedright\arraybackslash}p{0.3\textwidth}@{}}\toprule Field & Type & Req. & Unit & Range\\\midrule
\texttt{\scriptsize schema} & const & yes & - & = example.readiness/v1\\
\texttt{\scriptsize product\_version} & string & yes & - & pattern\\
\texttt{\scriptsize catalog\_sha256} & string & yes & - & pattern\\
\texttt{\scriptsize contracts} & array of object & yes & - & 1 to items\\
\bottomrule\end{tabular}\end{table}


\subsection{Read-only API (api)}\label{sec:data-api}

The api contract (version \cver{api}) is the dormant read-only layer over the committed documents: GET only, with the index, a case manifest, an artifact by path, the benchmark results, the model metrics and the readiness document (Table~\ref{tab:contract-api}). It is owned by the application folder, tested against the OpenAPI document of the catalog, and consumed by the typed client of the frontend. While no document is baked, or while a route's document is absent, the API answers 503 with the state missing and the reason, never an empty body; an unknown case or artifact is 404. The static readiness says missing for this row while the API lane is dormant, and the dynamic readiness route says available. The API is the only boundary whose producer is a service and not a file, which is why its contract is the OpenAPI document and not a schema.

\begin{table}[t]\centering\footnotesize
\caption{Operations of the dormant read-only API contract (the OpenAPI document of the catalog): method and path, the operation and its response codes; every other method is refused.}\label{tab:contract-api}
\begin{tabular}{@{}l >{\raggedright\arraybackslash}p{0.4\textwidth} l@{}}\toprule Operation & Summary & Responses\\\midrule
\texttt{\scriptsize GET /health} & Liveness. & 200\\
\texttt{\scriptsize GET /healthz} & Liveness (alias of /health). & 200\\
\texttt{\scriptsize GET /api/cases} & The case index (catalog row case-index), baked and refused cases. & 200, 503\\
\texttt{\scriptsize GET /api/cases/\{case\_id\}/manifest} & One case manifest (catalog row case-manifest). & 200, 404, 503\\
\texttt{\scriptsize GET /api/artifacts/\{rel\_path\}} & An artifact a manifest names (catalog row trace), by its path under data/derived. & 200, 404, 503\\
\texttt{\scriptsize GET /api/benchmark} & The benchmark results (catalog row benchmark-results). & 200, 503\\
\texttt{\scriptsize GET /api/models/metrics} & The held-out metrics of the model slots (catalog row model-metrics). & 200, 503\\
\texttt{\scriptsize GET /api/readiness} & What is real in this build (catalog row readiness). & 200, 503\\
\bottomrule\end{tabular}\end{table}

